% titel: Blatt 3, A5: Näherung der Reihe 1/k² mit Python

% Teil I: Werkzeuge, kein Lösungsweg.

% Grundidee: sum_{k=1}^inf 1/k^2 = pi^2/6 (Euler). Man nähert die
% Teilsumme mit einer Schleife an und sucht das kleinste N, ab dem die
% Differenz < 0.001 ist. Typische Python-Fallen: range halb-offen, **.

Diese Aufgabe verbindet Analysis (Euler-Reihe) mit Programmierung. Die
Reihe $\sum 1/k^2$ konvergiert zu $\frac{\pi^2}6$, aber recht \emph{langsam}
--- das Restglied $=\sum_{k>N}1/k^2$ klingt wie $1/N$ ab. Genau deshalb ist
das gesuchte $N$ f\"ur die Schranke $0.001$ deutlich \emph{gr\"o\ss}er als
``bei den ersten Paar Gliedern``.

\begin{werkzeugkasten}[Python-Werkzeuge]
\begin{itemize}
  \item \texttt{math.pi} ist $\pi$, \texttt{import math} zuerst.
  \item \texttt{range(1,101)} liefert $1,2,\dots,100$ ($1$ bis $100$,
        oberer Wert exklusiv). Teilsumme bis $N$:
        \texttt{sum([1/k**2 for k in range(1, N+1)])}.
  \item Differenz: \texttt{math.pi**2/6 - Teilsumme}.
\end{itemize}
\end{werkzeugkasten}

\begin{werkzeugkasten}[Euler-Reihe und Restglied]
\begin{equation}
  \sum_{k=1}^{\infty}\frac1{k^2}=\frac{\pi^2}6 .
\end{equation}
Sie konvergiert (das \emph{Restglied} geht "wie $\int_N^\infty
\frac{dx}{x^2}=\frac1N$`` zu $0$). Um ``erstmals unter $0.001$`` zu
finden, erh\"oht man $N$ so lange, bis die Differenz kleiner als
$0.001$ wird.
\end{werkzeugkasten}

\begin{achtung}[Typische Fehler]
\begin{itemize}
  \item \texttt{range(1,N+1)} --- \texttt{range} endet nicht am oberen
        Wert.
  \item \texttt{k**2} statt \verb|k^2| (in Python ist \verb|^| ein
        Bit-XOR, kein Potenz-Operator).
  \item ``Erstmals unter $0.001$``: das kleinste $N$ mit Differenz
        $<0.001$, im aufsteigenden Lauf gefunden (bei $N=100$ ist die
        Differenz noch \emph{oberhalb} der Schranke).
\end{itemize}
\end{achtung}
